%*******************************************************
% Chapter 4
%*******************************************************
% STRUCTURE : brouillon 2, 23 avril 2026
% §4.0 Ce que cette thèse cherche à montrer (NEW)
%   §4.0.1 Proposition centrale
%   §4.0.2 Positionnement
%   §4.0.3 Hypothèses testables (H1-H4)
%   §4.0.4 Familles de scénarios
%   §4.0.5 Limites
%   §4.0.6 Plan d'expérience IMACLIM (tableau)
% §4.1 Le problème de la modélisation énergie-économie
% §4.2 Genèse et architecture d'IMACLIM
% §4.3 Le second-best
% §4.4 IMACLIM dans le paysage IAM
% §4.5 Pourquoi IMACLIM pour la question du numérique
%*******************************************************

\chapter{Design expérimental~: IMACLIM et représentation de l'ICT}\label{ch:imaclim}


\section{Ce que cette thèse cherche à montrer}\label{sec:ch4-hypotheses}


\subsection{Proposition centrale}\label{sec:ch4-these}

Le couplage entre les systèmes de l'information et les systèmes de l'énergie possède quatre propriétés économiques identifiables~: l'irréversibilité du capital investi, la faible sensibilité de la demande de calcul au coût de l'énergie, le décalage entre les horizons temporels des décisions d'investissement, et la déformation des mécanismes de coordination par les prix. Selon les conditions institutionnelles dans lesquelles il opère, ce couplage peut orienter la trajectoire du système énergétique vers la décarbonation ou vers l'intensification de la consommation.

Cette proposition est conditionnelle~: le même ensemble de propriétés produit des résultats opposés selon le contexte. La tâche du dispositif expérimental est de spécifier ces conditions \emph{avant} l'exécution des scénarios, de les traduire en paramètres du modèle, puis de vérifier si les résultats confirment ou contredisent l'orientation attendue\footnote{La liste de ces quatre propriétés est construite pour le besoin du dispositif IMACLIM et ne prétend pas à l'exhaustivité. D'autres propriétés --- effets de réseau, externalités de données, coûts de sortie des plateformes --- sont discutées dans la littérature mais ne sont pas formalisables dans un modèle d'équilibre général calculable.}.


\subsection{Positionnement}\label{sec:ch4-positionnement}

Deux travaux récents modélisent l'impact de la numérisation sur les trajectoires d'émissions dans des modèles d'optimisation à anticipations parfaites (TIAM et MESSAGEix)\footnote{Beath \emph{et al.} (2026, preprint Research Square, doi:10.21203/rs.3.rs-8582166/v1) et Wilson \emph{et al.} (2026, preprint Research Square, doi:10.21203/rs.3.rs-8941019/v1). Ce dernier papier utilise les projections de transformation digitale de Fan \emph{et al.} (2025, \emph{npj Climate Action} 4:79).}. Trois résultats en émergent.

Premièrement, l'incertitude sur les émissions en 2050 liée à la seule numérisation est du même ordre de grandeur que celle liée à l'ensemble des variables socioéconomiques et démographiques des scénarios SSP1-5. La numérisation n'est pas un facteur marginal~: c'est une variable de premier ordre pour les trajectoires climatiques.

Deuxièmement, les effets indirects de la numérisation --- les modifications de la demande de transport, de chauffage, de production industrielle par les applications numériques --- dominent les effets directs --- la consommation des centres de données --- d'un facteur cinq environ. Le transport passager est le secteur où l'écart entre scénarios favorables et défavorables est le plus grand.

Troisièmement, les évaluations existantes à court terme présentent un biais d'optimisme~: elles sélectionnent des applications à effet positif, définissent des périmètres qui incluent les gains d'efficacité mais excluent le rebond et la demande induite. Les modèles d'affaires actuels de l'IA générative --- orientés vers la consommation, le marketing et l'assistance personnelle --- sont caractéristiques des conditions qui produisent les résultats défavorables, pas les résultats favorables.

Ces travaux traitent la numérisation comme une perturbation appliquée de l'extérieur à un système énergétique par ailleurs inchangé. Le couplage y est un modificateur de demande sectorielle, pas une transformation de la structure du système. C'est l'écart que le dispositif proposé ici vise à combler~: modéliser le couplage dans un cadre où le capital est irréversible, les agents sont myopes, le marché du travail ne s'ajuste pas instantanément, et les distorsions de marché interagissent entre elles.

Le dispositif expérimental mobilise deux des six méthodes d'évaluation des modèles intégrés identifiées par \textcite{wilson_evaluating_2021}. Premièrement, l'analyse de sensibilité~: le balayage paramétrique des hypothèses H1 à H4 (irréversibilité, élasticité, rebond, fracture numérique) produit des enveloppes de résultats qui relient la structure du modèle à son comportement. Deuxièmement, la confrontation aux faits stylisés~: les régularités historiques établies dans les chapitres~1 et~2 (co-saturations transport-communication, gradient isolé/intégré du capital, périodisation du rapport improvement/énergie) constituent des contraintes que les trajectoires projetées par le modèle doivent respecter. Si une trajectoire projetée viole un fait stylisé établi, c'est un signal d'erreur dans le modèle, pas une prédiction. Les résultats publiés de MESSAGEix \parencite{wilson_impacts_2026} et de TIAM \parencite{beath_impacts_2026} sur le même objet fournissent des points de repère utiles pour interpréter les trajectoires IMACLIM~: les différences de résultats, quand elles apparaissent, sont attribuables aux différences de structure (anticipations parfaites vs.~myopiques, capital putty-putty vs.~putty-clay, couplage exogène vs.~endogène), non à des divergences sur les données d'entrée. Une comparaison inter-modèles formelle dépasse toutefois le périmètre de cette thèse. Cette évaluation vise à établir l'\emph{appropriateness} (le modèle répond-il à la bonne question~?), l'\emph{interpretability} (les résultats sont-ils compréhensibles à la lumière de la structure~?) et la \emph{credibility} (le modèle est-il perçu comme adéquat par ses utilisateurs~?) du dispositif.


\subsection{Hypothèses testables}\label{sec:ch4-soustheses}

Quatre hypothèses, directement testables dans le cadre IMACLIM, structurent le dispositif.

\paragraph{H1~--- L'irréversibilité du capital verrouille les trajectoires.}
L'infrastructure numérique lourde (centres de données, réseaux, équipements de refroidissement) a une durée de vie de 20 à 25~ans. Les centrales qui l'alimentent ont une durée de vie de 30 à 40~ans. Les modèles d'IA qu'elle sert ont une durée de vie de 12 à 18~mois. Ce décalage fait que les décisions d'investissement prises aujourd'hui, sous une incertitude radicale sur la demande future de calcul, déterminent la trajectoire d'émissions pour les décennies suivantes --- indépendamment des politiques climatiques introduites par la suite. Le test~: comparer les trajectoires d'émissions dans IMACLIM avec du capital rigide (putty-clay) et avec du capital malléable dans le module numérique, toutes choses égales par ailleurs.

\paragraph{H2~--- La demande de calcul ne réagit pas au coût de l'énergie comme les autres demandes.}
Les utilisateurs finaux de calcul (entreprises, ménages, administrations) perçoivent un prix du service qui ne reflète pas le coût énergétique réel de production. L'énergie représente une fraction modeste du coût total du service numérique pour l'utilisateur, exactement comme elle représentait une fraction négligeable du coût du transport ferroviaire au XIX\textsuperscript{e}~siècle ou du coût de la radiodiffusion dans les années~1920. Cette propriété implique que le signal de prix de l'énergie, même s'il fonctionne côté offre (les opérateurs de centres de données localisent leurs installations en fonction du coût de l'électricité), ne freine pas la croissance de la demande de calcul de la manière qu'un modèle standard prédirait\footnote{Nuance importante~: la faible élasticité concerne la \emph{demande de service} (les utilisateurs finaux), pas nécessairement les \emph{opérateurs d'infrastructure}, qui sont sensibles au coût de l'énergie dans leurs décisions de localisation et d'approvisionnement.}. Le test~: varier l'élasticité-prix de la demande de calcul et identifier le seuil au-delà duquel la trajectoire d'émissions change de régime qualitatif.

\paragraph{H3~--- Le rebond est plus fort dans un monde imparfait que dans un monde optimal.}
Quand une application numérique réduit le coût d'une activité (le coût du kilomètre parcouru, le coût de chauffage par mètre carré, le coût de production d'une unité industrielle), les agents réinvestissent une partie des gains en activité supplémentaire. C'est le paradoxe de Jevons. L'hypothèse testée ici est que ce rebond est structurellement plus fort dans un modèle avec agents myopes et capital irréversible que dans un modèle d'optimisation à anticipations parfaites, parce que les agents ne réallouent pas instantanément les gains vers les usages les plus productifs~: ils les consomment. Le test~: appliquer les mêmes perturbations de demande sectorielle dans IMACLIM et dans un modèle d'optimisation, et comparer les résultats.

\paragraph{H4~--- La fracture numérique crée une asymétrie géographique dans les trajectoires d'émissions.}
Les mécanismes du couplage ne s'activent pas uniformément dans le monde. Ils sont concentrés dans les régions à forte transformation digitale (Amérique du Nord, Europe, Asie de l'Est), où ils ajoutent de la demande au système énergétique. Dans les régions à faible transformation digitale, la numérisation n'a ni les effets positifs (optimisation du système) ni les effets négatifs (demande additionnelle) qu'on observe ailleurs. Le bilan net global dépend de cette géographie. Le test~: différencier l'intensité des perturbations par niveau de transformation digitale régionale\footnote{La difficulté méthodologique est d'isoler l'effet de la fracture numérique de celui de la stagnation économique générale, qui covarient dans les scénarios SSP3.}.


\subsection{Familles de scénarios}\label{sec:ch4-scenarios}

\wip{Cadrage préliminaire. Les paramètres de chaque scénario seront spécifiés en détail dans la section sur le dispositif expérimental.}

La question posée dans cette thèse --- sous quelles conditions le couplage ICT-énergie intensifie ou atténue la consommation --- n'appelle pas un scénario unique mais un espace de futurs contrastants. L'objectif des scénarios n'est pas de prédire lequel se réalisera~; il est de délimiter l'enveloppe des trajectoires possibles et d'identifier les paramètres structurels qui font basculer le résultat d'un régime à un autre. Un scénario est un monde économique cohérent, défini par des hypothèses sur la structure des marchés, les comportements des agents et les contraintes physiques~; ses conséquences climatiques sont calculées par le modèle, non présupposées par le modélisateur.

Cette conception des scénarios comme mondes structurels se distingue de trois approches alternatives. Les scénarios d'optimisation, tels que ceux produits par REMIND ou MESSAGE, calculent la trajectoire la moins coûteuse sous une contrainte climatique donnée~; ils répondent à la question \og que faut-il faire~?\fg{} et présupposent des agents omniscients et des marchés parfaits. Les scénarios narratifs binaires, tels que les scénarios \emph{Enable} et \emph{Undermine} de \textcite{wilson_impacts_2026}, distinguent un avenir favorable d'un avenir défavorable mais n'expliquent pas par quels mécanismes économiques l'un ou l'autre advient~; la section précédente a montré que le mot \og systémique\fg{} y recouvre des canaux hétérogènes que la catégorie unifiée empêche de discriminer (\S\ref{subsec:scenarisation}). Les scénarios normatifs, tels que ceux du WBGU ou de l'ADEME \emph{Transitions~2050}, définissent un futur souhaitable puis construisent le chemin pour l'atteindre~; ils répondent à la question \og que voulons-nous~?\fg{} et supposent que la direction de la transition est connue d'avance. L'approche retenue ici, inspirée de \textcite{guivarch_costs_2011}, répond à une question différente~: \og que se passe-t-il si la structure de l'économie est celle-ci plutôt que celle-là~?\fg{} Les scénarios ne sont ni optimaux, ni souhaitables, ni probables~; ils sont cohérents, et la cohérence est assurée par la discipline du modèle.

La construction de ces mondes procède en quatre opérations, qui correspondent à quatre traditions méthodologiques et ne doivent pas être confondues.

La première opération est la définition des mondes eux-mêmes. Chaque scénario est une configuration d'hypothèses sur les paramètres du modèle~: taux de marge, élasticités, taux de dépréciation, contraintes de capacité, coefficients techniques. Ces paramètres ne sont pas des variables libres~; ils sont adossés aux mécanismes économiques identifiés dans le chapitre~2 et aux faits stylisés documentés dans les chapitres~1 et~3. La méthode est celle de \textcite{guivarch_costs_2011}~: les scénarios sont des mondes seconds, définis par des structures de marché imparfaites et des comportements d'agents myopes. La différence entre deux scénarios est une différence de structure, non de préférence normative. C'est cette première opération qui distingue le dispositif de la présente thèse de celui de \textcite{wilson_impacts_2026} et de \textcite{beath_impacts_2026}~: là où ces travaux appliquent des perturbations de demande sectorielle à un modèle dont la structure reste inchangée, les scénarios IMACLIM modifient la structure elle-même --- le régime du capital, la distribution des revenus, le fonctionnement du signal de prix.

La deuxième opération est la traduction des narratifs en paramètres. Un monde structurel ne naît pas d'un exercice de sensibilité~; il naît d'un récit économique qui identifie quels mécanismes du chapitre~2 sont actifs, à quelle intensité, et comment ils interagissent. \textcite{wilson_impacts_2026} ont codifié cette traduction pour la numérisation dans MESSAGEix, en articulant choix d'impacts par application, estimation de la fourchette empirique, règle de diffusion géographique et temporelle, et contexte d'interprétation. Le même geste est reproduit ici, à la différence que les mécanismes traduits ne se limitent pas aux modificateurs de demande sectorielle~: ils incluent l'irréversibilité du capital, les taux de marge, la distribution fonctionnelle des revenus et la déformation du signal de prix. Pour chaque dimension du chapitre~2, la question posée est~: dans ce monde, quelle est l'hypothèse active~? Le tableau dimension par dimension qui en résulte constitue la spécification complète du scénario avant toute exécution du modèle.

La troisième opération est la vérification de la couverture de l'espace des possibles par analyse morphologique. Cette méthode, formalisée par Zwicky puis largement diffusée en prospective française par \textcite{godet_creating_2006} et l'institut Futuribles, consiste à décomposer un problème en variables structurantes, à formuler pour chacune plusieurs hypothèses d'évolution mutuellement exclusives, puis à engendrer des scénarios en combinant une hypothèse par variable. La méthode ne présuppose pas de théorie causale~; elle explore systématiquement l'espace combinatoire pour repérer des configurations que le raisonnement par traditions intellectuelles n'aurait pas envisagées.

\textcite{deron_par_2025}, dans le cadre du projet Chemins de transition (Université de Montréal), appliquent cette méthode à la question du numérique dans un monde contraint par les limites planétaires. Leur tableau morphologique retient cinq variables~: l'état de la géopolitique mondiale (religions comme source de légitimation, intégration continentale, fin de l'OTAN, réparation solidaire), le modèle de gouvernance locale (décentralisée, jurys citoyens, étatique descendante, vote direct), le modèle de gestion du numérique (communautaire, surveillance, service public, paiement direct), le contexte de production des appareils (spécialisation régionale, fabrication locale, chaîne mondialisée automatisée, circularité), et la place du numérique dans la société (accès inégal, combinaison low-tech et high-tech, usages collectivisés, ubiquité sans obsolescence). Le tableau a été soumis à 25~personnes issues de différentes disciplines et enrichi d'une centaine de commentaires, ce qui a permis d'intégrer une diversité épistémique que les quatre auteurs, par leur proximité disciplinaire, n'auraient pas produite seuls. Les hypothèses combinent des évolutions tendancielles --- prolongation des dynamiques observées --- et des évolutions exploratoires fondées sur des signaux faibles, ce qui permet d'explorer des régions du cône des futurs éloignées des extrapolations habituelles.

Quatre scénarios narratifs résultent de cette procédure. Leur intérêt pour la présente thèse tient moins aux récits eux-mêmes qu'à trois propriétés de la méthode. Premièrement, elle force des combinaisons non évidentes~: une surveillance généralisée dans une société low-tech, un numérique circulaire dans un régime de réparation solidaire internationale. Ces configurations montrent que les scénarios les plus réalistes ne sont pas nécessairement les plus purs. Deuxièmement, elle rend visibles les variables que l'on oublie quand on raisonne par traditions intellectuelles~; en l'occurrence, le contexte de production des appareils (chaînes d'approvisionnement, relocalisation, circularité) ne figure dans aucune des dix dimensions du chapitre~2, ce qui constitue une limite explicite du dispositif IMACLIM. Troisièmement, la mise en récit oblige à vérifier la cohérence interne de chaque combinaison, exercice que le paramétrage d'un modèle ne garantit pas à lui seul.

La procédure morphologique ne produit pas les scénarios de la thèse~; elle les teste. Les scénarios ont été construits par la première opération, \`a partir de traditions intellectuelles et des mécanismes du chapitre~2. L'analyse morphologique intervient ensuite pour vérifier leur cohérence, leur couverture de l'espace des possibles et leur non-redondance. Elle est un outil de contr\^ole, non un générateur automatique de futurs.

Le protocole proc\`ede en trois temps~: le choix des variables structurantes, la construction du tableau morphologique, et la projection des scénarios existants sur ce tableau.

\emph{Choix des variables.} Les variables morphologiques ne sont ni les scénarios, ni les hypoth\`eses du chapitre~2 prises une \`a une, ni le couplage ICT-\'energie lui-m\^eme. Elles doivent \^etre les conditions structurantes qui font varier les scénarios et qui permettent ensuite d'observer quel type de couplage \'emerge. Quatre crit\`eres guident leur s\'election.

Le premier crit\`ere est la \emph{discrimination}~: la variable doit distinguer au moins trois sc\'enarios. Si tous les sc\'enarios prennent la m\^eme valeur sur une variable, celle-ci ne structure pas l'espace et doit \^etre \'ecart\'ee. Le deuxi\`eme crit\`ere est l'\emph{ext\'eriorit\'e relative}~: la variable ne doit pas pr\'esupposer le r\'esultat que l'on cherche \`a observer. Le \og couplage fort\fg{} ou l'\og internalisation des co\^uts\fg{} ne peuvent donc pas \^etre des variables~; ce sont des r\'esultats. Le troisi\`eme crit\`ere est l'\emph{ancrage th\'eorique}~: la variable doit \^etre reli\'ee aux dimensions du chapitre~2 (D1--D10) ou aux transversales des chapitres~1 et~3 (\'energie-mati\`ere, gouvernance), de sorte que chaque hypoth\`ese morphologique puisse \^etre renvoy\'ee \`a un m\'ecanisme \'economique identifi\'e. Le quatri\`eme crit\`ere est la \emph{traductibilit\'e}~: la variable doit pouvoir, au moins partiellement, devenir un param\`etre, une contrainte ou une hypoth\`ese narrative dans IMACLIM. Une variable purement sociologique ou culturelle, si elle n'a pas de traduction param\'etrique, peut figurer dans le tableau mais ne produira pas de test mod\'elisable.

\emph{Variables retenues.} Cinq variables satisfont ces quatre crit\`eres. Elles sont organis\'ees du contexte le plus large au m\'ecanisme le plus sp\'ecifique, dans un ordre qui fait \'echo \`a la structure de Deron et al.\ (deux variables de contexte g\'en\'eral, trois variables sp\'ecifiques \`a l'objet d'\'etude) tout en restant ancr\'e dans l'architecture des chapitres pr\'ec\'edents.

La premi\`ere variable, \emph{ordre mondial et cha\^ines de valeur}, d\'ecrit l'\'echelle \`a laquelle les syst\`emes d'information et d'\'energie sont organis\'es~: mondialisation int\'egr\'ee, blocs rivaux, relocalisation, ou r\'eparation solidaire. Elle est ancr\'ee dans le chapitre~1 (les phases historiques de mondialisation et de fragmentation des r\'eseaux), dans la dimension D9 (g\'eographie) et dans la transversale gouvernance. Elle se traduit dans IMACLIM par la diff\'erenciation r\'egionale des param\`etres (12~r\'egions).

La deuxi\`eme variable, \emph{r\'egime du capital informationnel}, d\'ecrit qui construit, poss\`ede, finance et verrouille l'infrastructure de calcul~: accumulation concentr\'ee (hyperscale, int\'egr\'ee verticalement, autofinanc\'ee par les rentes), capital distribu\'e et communs (coop\'eratives, open source, infrastructure l\'eg\`ere), capital r\'eorient\'e (stock existant redirig\'e apr\`es une crise), ou capital contraint par raret\'e (goulots amont sur les semi-conducteurs, les m\'etaux ou l'\'energie). Elle est ancr\'ee dans les dimensions D2 (capital), D4 (rendements et r\'eseaux) et D8 (financement), et enrichie par les cas historiques du chapitre~1 (Edison, Bell System, IBM, hyperscalers). Elle se traduit dans IMACLIM par les taux de marge~($\pi_j$), les taux de d\'epr\'eciation~($\delta_j$) et les profils d'investissement.

La troisi\`eme variable, \emph{direction du changement technique}, d\'ecrit vers quoi l'effort d'innovation s'oriente~: puissance brute du calcul (\emph{compute-first}), efficacit\'e \'energ\'etique (\emph{energy-aware}), ou appropriabilit\'e et suffisance (robustesse, r\'eparabilit\'e, ad\'equation \`a l'usage). Elle est ancr\'ee dans la dimension D10 (innovation) et la dimension D3 (travail), car la direction du changement technique d\'etermine le rapport travail-automation. Le chapitre~1 montre que chaque phase historique des syst\`emes d'information a eu sa direction technique dominante~; le chapitre~3 y ajoute la \emph{quality-energy ladder} qui relie le progr\`es technique \`a l'intensit\'e \'energ\'etique. Cette variable se traduit dans IMACLIM par les coefficients techniques et les trajectoires d'efficacit\'e dans les modules dynamiques.

La quatri\`eme variable, \emph{dynamiques \'energ\'etiques et mat\'erielles}, d\'ecrit l'\'etat et l'\'evolution du syst\`eme \'energ\'etique et des ressources mat\'erielles qui portent l'infrastructure num\'erique~: abondance d\'ecarbon\'ee (d\'eploiement rapide des renouvelables, mat\'eriaux disponibles), transition lente (fossiles r\'esilients, changement incr\'emental), contrainte mat\'erielle (\'energie, m\'etaux, eau, sites en raret\'e g\'en\'eralis\'ee), ou abondance carbon\'ee (\'energie fossile bon march\'e, dommages climatiques diff\'er\'es). Elle est ancr\'ee dans la transversale \'energie-mati\`ere (0-TER) et dans le chapitre~3 qui en \'etablit les fondements biophysiques. Elle se traduit dans IMACLIM par la trajectoire du mix \'electrique, les prix des \'energies primaires et les contraintes de capacit\'e.

La cinqui\`eme variable, \emph{nature et usages de la demande de calcul}, d\'ecrit pourquoi et \`a quelle sensibilit\'e au co\^ut le calcul est demand\'e~: demande induite par l'offre avec co\^ut \'energ\'etique masqu\'e (plateformes, mod\`ele prix-z\'ero, contrats d'achat direct), demande tir\'ee par la productivit\'e (automation B2B, gains d'efficacit\'e), demande born\'ee par les pratiques (culture, suffisance, discernement), demande contrainte par la raret\'e physique, ou demande disciplin\'ee par l'allocation politique (quotas, rationnement, priorisation sectorielle). Cette variable absorbe les dimensions D5 (demande), D6 (signal de prix) et D4 (rendements de r\'eseau) du chapitre~2. C'est elle qui s\'epare des sc\'enarios qui seraient autrement indistinguables sur les autres axes~: S1 (demande induite par la rente) et S3 (demande tir\'ee par la productivit\'e) partagent le m\^eme r\'egime de capital mais diff\`erent par le moteur de la demande. Elle se traduit dans IMACLIM par les \'elasticit\'es-prix, les \'elasticit\'es-revenu et les trajectoires exog\`enes de demande sectorielle.

Aucune de ces cinq variables ne d\'ecrit directement le couplage ICT-\'energie~; elles d\'ecrivent les conditions sous lesquelles il prend forme. Le couplage est le r\'esultat de la combinaison, non une hypoth\`ese impos\'ee au d\'epart.

Le tableau~\ref{tab:morpho} pr\'esente la grille morphologique compl\`ete.

\begin{table}[htbp]
\centering
\small
\caption{Tableau morphologique~: cinq variables structurantes et leurs hypoth\`eses d'\'evolution. T~= tendancielle, E~= exploratoire.}\label{tab:morpho}
\begin{tabular}{p{3.8cm}p{2.2cm}p{2.2cm}p{2.2cm}p{2.2cm}p{2.2cm}}
\hline
\textbf{Variable} & \textbf{H1} & \textbf{H2} & \textbf{H3} & \textbf{H4} & \textbf{H5} \\
\hline
V1~--- Ordre mondial et cha\^ines de valeur &
Mondialisation int\'egr\'ee~(T) &
Blocs rivaux~(T/E) &
Relocalisation~(E) &
R\'eparation solidaire~(E) & \\
\hline
V2~--- R\'egime du capital informationnel &
Accumulation concentr\'ee~(T) &
Distribu\'e --- communs~(E) &
R\'eorient\'e~(E) &
Contraint par raret\'e~(E) & \\
\hline
V3~--- Direction du changement technique &
\emph{Compute-first}~(T) &
\emph{Energy-aware}~(T/E) &
Appropriabilit\'e --- suffisance~(E) & & \\
\hline
V4~--- Dynamiques \'energ\'etiques et mat\'erielles &
Abondance d\'ecarbon\'ee~(E) &
Transition lente~(T) &
Contrainte mat\'erielle~(E) &
Abondance carbon\'ee~(T) & \\
\hline
V5~--- Nature et usages de la demande &
Induite par l'offre, co\^ut masqu\'e~(T) &
Tir\'ee par la productivit\'e~(T) &
Born\'ee par les pratiques~(E) &
Contrainte par raret\'e~(T/E) &
Allocation politique~(E) \\
\hline
\end{tabular}
\end{table}

\emph{Projection des sc\'enarios.} Les mondes structurels d\'efinis par la premi\`ere op\'eration sont projet\'es sur la grille en cochant, pour chaque sc\'enario, une hypoth\`ese par variable. Le tableau~\ref{tab:projection} pr\'esente le r\'esultat. Trois observations en ressortent. Premi\`erement, chaque paire de sc\'enarios diff\`ere sur au moins deux variables, ce qui confirme leur non-redondance. Deuxi\`emement, toutes les hypoth\`eses apparaissent au moins une fois, \`a l'exception de \og r\'eparation solidaire\fg{} en~V1, qui sort du p\'erim\`etre du dispositif IMACLIM. Troisi\`emement, S4 est le seul sc\'enario dont le profil morphologique change en cours de route~: avant le \emph{turning point}, il occupe les m\^emes cases que~S1~; apr\`es, il bascule sur quatre variables sur cinq. La relation S4-sans-\emph{turning-point}~$=$~S1 est lisible dans le tableau et testable dans le mod\`ele.

\begin{table}[htbp]
\centering
\small
\caption{Projection des sc\'enarios sur le tableau morphologique.}\label{tab:projection}
\begin{tabular}{p{2.8cm}p{2cm}p{2.2cm}p{2cm}p{2cm}p{2.5cm}}
\hline
& \textbf{V1} & \textbf{V2} & \textbf{V3} & \textbf{V4} & \textbf{V5} \\
\hline
S0b Prolongation & Int\'egr\'ee & Concentr\'e & Compute-first & Trans.\ lente & Induite, masqu\'e \\
S1 Accumulation & Int\'egr\'ee & Concentr\'e & Compute-first & Trans.\ lente & Induite, masqu\'e \\
S2 Substrat & Int\'egr\'ee & Contraint & Energy-aware & Contrainte & Raret\'e physique \\
S3 Productivit\'e & Int\'egr\'ee & Concentr\'e & Energy-aware & Trans.\ lente & Productivit\'e \\
S4 avant TP & Int\'egr\'ee & Concentr\'e & Compute-first & Trans.\ lente & Induite, masqu\'e \\
S4 apr\`es TP & Int\'egr\'ee & R\'eorient\'e & Energy-aware & D\'ecarbon\'ee & Productivit\'e \\
S5 Appropriation & Relocalis\'ee & Communs & Appropriabilit\'e & Trans.\ lente & Pratiques \\
S6 Fragmentation & Blocs rivaux & Concentr\'e & (par bloc) & (par bloc) & (par bloc) \\
S7 Raret\'e gouv. & Relocalis\'ee & Contraint & Energy-aware & Contrainte & Alloc.\ politique \\
\hline
\end{tabular}
\end{table}

L'exercice a \'egalement identifi\'e des combinaisons non triviales que les sc\'enarios retenus ne couvrent pas~: un monde o\`u le compute est intense et d\'ecentralis\'e (\emph{compute-first} sous relocalisation), un monde o\`u les communs num\'eriques coexistent avec une automation substitutive forte, ou encore un r\'egime d'accumulation concentr\'ee sous d\'ecarbonation rapide et r\'egulation institutionnelle (entre~S1 et~S4). Ces combinaisons enrichissent la discussion du chapitre~\ref{ch:discussion} sans constituer des sc\'enarios IMACLIM \`a part enti\`ere.

La quatrième opération est le positionnement dans le cadre de référence du champ. Les scénarios de cette thèse doivent pouvoir être situés par rapport aux trajectoires socioéconomiques partagées (SSP) qui structurent la modélisation climatique contemporaine, ne serait-ce que parce que les projections de \textcite{fan_digital_2025} fournissent les niveaux de transformation numérique régionale sous chaque SSP. Les mondes structurels ne sont pas des sous-variantes d'un SSP particulier~; ils croisent plusieurs SSP selon les dimensions considérées. Le test de cohérence est de vérifier que chaque monde est compatible avec au moins une trajectoire SSP sur les variables qu'il ne spécifie pas directement.

Ces quatre opérations --- mondes structurels, traduction paramétrique, vérification combinatoire, positionnement SSP --- produisent les familles de scénarios qui suivent. Les scénarios ne sont pas des récits sur l'avenir du monde. Ce sont des configurations de paramètres qui activent ou désactivent des mécanismes, à des intensités différentes, pour tester les hypothèses formulées ci-dessus. Les familles définies ci-dessous portent sur le côté énergie-physique du couplage. Chacune peut être déclinée selon un axe macroéconomique orthogonal~--- la direction du changement technique numérique au sens d'\textcite{acemoglu_building_2026}~--- qui se traduit par la trajectoire de la part du travail dans la valeur ajoutée (variante \og automation-dominante\fg{} vs.\ variante \og pro-worker\fg{}). Ce croisement est détaillé dans le plan d'expérience (tableau~\ref{tab:plan-imaclim}).

\paragraph{S0a~--- Contrefactuel (mod\`ele aveugle).}
L'ICT est un secteur standard sans propri\'et\'es sp\'ecifiques dans IMACLIM~: ni irr\'eversibilit\'e diff\'erenci\'ee, ni \'elasticit\'e-prix r\'eduite, ni perturbation sectorielle. Ce contrefactuel mesure ce que l'endog\'en\'eisation du couplage change par rapport \`a un monde o\`u le num\'erique n'a pas de propri\'et\'es distinctives.

\paragraph{S0b~--- Prolongation.}
Les structures actuelles se prolongent. L'accumulation du capital num\'erique continue \`a son rythme tendanciel, la direction du changement technique reste orient\'ee vers la puissance brute, le syst\`eme \'energ\'etique \'evolue lentement, et le co\^ut \'energ\'etique du calcul reste invisible pour l'utilisateur final. Pas de crise, pas de bifurcation, pas de politique transformatrice. Ce sc\'enario, ancr\'e dans la th\'eorie de la d\'ependance de sentier \parencite{arthur_competing_1989, david_clio_1985}, sert de baseline~: tous les autres sc\'enarios sont \'evalu\'es par \'ecart \`a celui-ci.

\paragraph{S1~--- Accumulation computationnelle.}
Le calcul devient le lieu principal d'accumulation du capital. Quelques firmes hyperscale concentrent l'infrastructure, l'autofinancent par les rentes, int\`egrent verticalement l'\'energie par des contrats d'achat direct, et cr\'eent leur propre demande par la boucle donn\'ees-qualit\'e-utilisateurs. La part du travail dans la valeur ajout\'ee baisse. Le co\^ut \'energ\'etique est absorb\'e dans la cha\^ine de valeur et ne parvient pas \`a l'utilisateur final. Ce sc\'enario est adoss\'e \`a la tradition marxiste et \`a l'analyse du capitalisme de plateforme \parencite{durand_techno-feodalisme_2020}~: le capital informationnel y fonctionne comme un r\'egime d'accumulation dont les contradictions --- suraccumulation, chute de la part du travail, \'epuisement du substrat mat\'eriel --- d\'eterminent les sorties possibles vers les autres sc\'enarios. Ce sc\'enario teste H1 (irr\'eversibilit\'e) et H2 (\'elasticit\'e-prix).

\paragraph{S2~--- Contraintes du substrat.}
Les limites mat\'erielles et \'energ\'etiques s'imposent au syst\`eme num\'erique. L'\'energie se rar\'efie ou rench\'erit, les m\'etaux critiques deviennent des goulots, l'eau manque, les sites sont satur\'es. Le capital num\'erique existant se d\'epr\'ecie non par usure mais par non-viabilit\'e \'energ\'etique --- ce que \textcite{deron_sprouting_2025} appellent des \og ruines\fg{} au sens de \textcite{monnin_communs_2021}. Le signal de prix est restaur\'e par la physique, pas par la politique. La direction du changement technique est forc\'ee vers l'\'economie d'\'energie, non par choix mais par n\'ecessit\'e. Ce sc\'enario est ancr\'e dans la tradition ricardienne et dans l'\'economie \'ecologique \parencite{georgescu-roegen_entropy_1971}. Il teste H1 (le capital verrouill\'e devient un actif \'echou\'e) et H4 (la contrainte est g\'eographiquement diff\'erenci\'ee).

\paragraph{S3~--- Productivit\'e fossile.}
Le changement technique est nominalement orient\'e vers l'efficacit\'e \'energ\'etique --- l'intensit\'e du calcul par watt progresse. Mais les gains d'efficacit\'e lib\`erent du pouvoir d'achat et des capacit\'es qui se r\'einvestissent en volume. La productivit\'e totale des facteurs est \'elev\'ee, l'intensit\'e baisse, mais la consommation totale monte. Le syst\`eme fossile est optimis\'e, pas remplac\'e. C'est le sc\'enario le plus d\'elicat \`a identifier car il ressemble \`a une transition --- les indicateurs d'intensit\'e s'am\'eliorent pendant que les \'emissions absolues augmentent. \textcite{deron_par_2025} observent que le discours du \og num\'erique responsable\fg{} --- \'eco-conception, mesure d'empreinte, sobri\'et\'e d'usage --- fonctionne localement mais ne modifie pas les dynamiques structurelles, ce qui en fait l'habillage id\'eologique de ce sc\'enario. Ancr\'e dans le paradoxe de Jevons et la litt\'erature sur le rebond \parencite{brockway_energy_2021}, ce sc\'enario teste H3 (rebond en monde imparfait) et H2 (signal de prix insuffisant).

\paragraph{S4~--- Cycle de d\'eploiement.}
Le d\'eploiement de la GPT num\'erique suit le sch\'ema identifi\'e par \textcite{perez_technological_2002} en deux phases. La phase d'installation, avant le \emph{turning point}, est indistinguable de~S1~: fr\'en\'esie sp\'eculative, accumulation concentr\'ee, direction technique orient\'ee vers la puissance brute, co\^ut \'energ\'etique masqu\'e. Le \emph{turning point} --- s'il advient --- est une crise financi\`ere ou institutionnelle qui redirige le capital accumul\'e vers des usages productifs et coordonn\'es. Apr\`es le \emph{turning point}~: le capital est r\'eorient\'e, la direction technique bascule vers l'efficacit\'e, le syst\`eme \'energ\'etique acc\'el\`ere sa d\'ecarbonation, et la demande est tir\'ee par la productivit\'e industrielle plut\^ot que par la rente. Ce sc\'enario a une propri\'et\'e structurelle unique~: il est le seul dont le profil morphologique change en cours de route. La relation S4-sans-\emph{turning-point}~$=$~S1 est testable dans IMACLIM~: si les conditions institutionnelles du \emph{turning point} ne sont pas r\'eunies, la fr\'en\'esie ne se r\'esout pas en d\'eploiement mais en accumulation permanente. Ce sc\'enario teste les quatre hypoth\`eses simultan\'ement.

\paragraph{S5~--- Appropriation technique.}
La demande de calcul est born\'ee par les pratiques, pas par le prix ni par la physique. Les usagers ne demandent pas davantage de puissance de calcul --- par discernement, par culture, par normes sociales. Le capital est distribu\'e, l\'eger, r\'eparable, interop\'erable. L'innovation vise l'appropriabilit\'e et la suffisance~: robustesse, r\'eparabilit\'e, ad\'equation \`a l'usage --- ce que la communaut\'e du \emph{permacomputing} et les travaux de \textcite{deron_sprouting_2025} d\'esignent comme les \og graines\fg{} d'un autre r\'egime technique. Ancr\'e dans la philosophie de la technique \parencite{stiegler_technique_1994, illich_tools_1973} et la th\'eorie des communs \parencite{ostrom_governing_1990}, c'est le seul sc\'enario dont le m\'ecanisme central est un changement de pr\'ef\'erences plut\^ot qu'un changement de structure. La question testable dans IMACLIM est de savoir si une demande born\'ee modifie significativement la trajectoire macro\'economique, ou si S5 reste sympathique mais marginal.

\paragraph{S6~--- Fragmentation.}
Le monde se scinde en blocs techniques et \'energ\'etiques. Chaque bloc r\'eplique ses cha\^ines de semi-conducteurs, ses infrastructures de calcul, ses standards. Les \'economies d'\'echelle mondiales sont perdues. Les co\^uts de production montent. Les trajectoires \'energ\'etiques divergent par bloc~: un bloc peut d\'ecarboner rapidement pendant qu'un autre reste fossile. Ancr\'e dans le r\'ealisme en relations internationales et la th\'eorie des cycles h\'eg\'emoniques \parencite{arrighi_long_1994}, ce sc\'enario ne prescrit que la variable g\'eopolitique~; le contenu interne de chaque bloc d\'epend du bloc. Dans IMACLIM-R World (12~r\'egions), il se traduit par des param\`etres diff\'erenci\'es par r\'egion. Ce sc\'enario teste H4 (fracture num\'erique et asym\'etrie g\'eographique).

\paragraph{S7~--- Raret\'e gouvern\'ee.}
L'\'energie et les mat\'eriaux ne suffisent pas --- comme dans S2 --- mais au lieu de laisser le march\'e rationner par les prix, l'autorit\'e publique d\'ecide qui a le droit de calculer. Quotas par secteur, priorit\'es d'usage, allocation administrative. Le signal de prix est remplac\'e par une d\'ecision politique. Ancr\'e dans l'\'economie du rationnement et de la planification \parencite{pigou_political_1920}, ce sc\'enario suppose \`a la fois la contrainte physique (comme~S2) et la volont\'e politique de la g\'erer (ce que~S2 ne suppose pas). La distinction S2/S7 est pr\'ecis\'ement la diff\'erence entre laisser le march\'e rationner et d\'ecider collectivement qui obtient quoi.

\medskip

Ces huit sc\'enarios ne sont pas ind\'ependants. S1 repr\'esente le r\'egime actuel pouss\'e \`a sa logique~; les autres sc\'enarios sont des r\'esolutions diff\'erentes de ses contradictions~: mat\'erielles (S2), institutionnelles (S4), sociales (S7), g\'eopolitiques (S6), culturelles (S5). S3 n'est pas une sortie de~S1 mais un S1 d\'eguis\'e en transition. S0b est un S1 att\'enu\'e. La structure dynamique du dispositif est donc~: S1 comme point de d\'epart analytique, les six autres comme modes de r\'esolution, et le mod\`ele comme arbitre de leurs cons\'equences climatiques.


\subsection{Ce que le dispositif ne couvre pas}\label{sec:ch4-limites}

Trois catégories d'effets sont exclues du périmètre. Les effets de la numérisation sur les institutions de gouvernance climatique (désinformation, polarisation, affaiblissement de la capacité régulatrice) ne sont pas formalisables dans un modèle d'équilibre général~; ils relèvent du narratif, pas du paramétrage. Les effets sur la productivité et l'emploi agrégés sont partiellement captés par la wage curve et le chômage endogène d'IMACLIM, mais la spécification du choc de productivité sectoriel lié à l'IA reste un problème ouvert~: les estimations publiées varient d'un facteur dix selon les auteurs et les méthodes. Les effets sur l'usage des sols (agriculture de précision, chaînes logistiques alimentaires) sont hors du périmètre de la configuration sectorielle actuelle d'IMACLIM (12~secteurs), non par choix conceptuel mais par contrainte technique.


\subsection{Plan d'expérience~: traduction dans IMACLIM}\label{sec:ch4-plan}

\wip{Tableau préliminaire. À affiner avec l'équipe IMACLIM du CIRED.}

Le tableau~\ref{tab:plan-imaclim} identifie, pour chaque hypothèse, ce qu'il faut modifier dans IMACLIM, quel module dynamique est concerné, quelles données sont nécessaires, et quelles sorties du modèle constituent le test.

\begin{table}[htbp]
\centering
\small
\caption{Plan d'expérience~: traduction des hypothèses dans IMACLIM-R World.}\label{tab:plan-imaclim}
\begin{tabular}{p{1.2cm}p{3cm}p{3.5cm}p{3cm}p{3cm}}
\hline
\textbf{Hyp.} & \textbf{Ce qu'on modifie} & \textbf{Où dans IMACLIM} & \textbf{Données nécessaires} & \textbf{Sortie testée} \\
\hline

H1 &
Durée de vie et rigidité du capital dans le secteur numérique &
\textbf{Modules dynamiques~:} investissement en capacité dans le secteur <<~services et industrie légère~>>. Paramètre de dépréciation du capital ($\delta_j$) et contrainte putty-clay. Comparer un run avec capital putty-clay (baseline IMACLIM) et un run avec capital putty-putty (modification) pour le sous-secteur numérique. &
Durée de vie des DC (20-25 ans), taux de renouvellement GPU (3-5 ans), profil CAPEX hyperscalers (IEA 2026~: 400 Md\$/an, +75\% en 2026). &
Trajectoire d'émissions, année de peak, sensibilité au timing de l'introduction d'une taxe carbone (avant vs.\ après 2035). \\

\hline

H2 &
Élasticité-prix de la demande d'électricité du secteur numérique &
\textbf{Équilibre statique~:} fonction de demande des ménages et des entreprises pour le bien <<~services~>>. Modifier l'élasticité-prix de la composante numérique (paramètre $\alpha_k$ dans la fonction d'utilité, ou coefficient technique $ic_{elec,services}$). Variantes~: élasticité standard ($-0{,}3$ à $-0{,}5$), élasticité réduite ($-0{,}05$ à $-0{,}1$), élasticité nulle. &
Axenbeck \& Niebel 2024 (firmes numérisées, Allemagne). IEA~: efficacité/tâche $\times$10/an \emph{mais} conso DC +50\%/an. &
Demande d'électricité agrégée, prix de l'électricité à l'équilibre, emploi sectoriel. Identifier le seuil d'élasticité qui change le régime qualitatif. \\

\hline

H3 &
Perturbations de demande sectorielle (<<~demand modifiers~>>) dans le transport et le bâtiment &
\textbf{Modules dynamiques~:} transport (taux de motorisation $\alpha_k$, distance parcourue, occupancy rate) et résidentiel/tertiaire (surface par tête, intensité énergétique). Appliquer les \emph{demand modifiers} Enable/Undermine calibrés par Wilson~\emph{et al.}~: transport $-20\%$ à $+81\%$, bâtiment $-22\%$ à $+9\%$. &
Wilson \emph{et al.} 2026 (Figure~6, uncertainty ranges par application). Rémi Paccou est co-auteur~: accès direct aux données. &
Comparer le rebond agrégé (émissions cumulées) avec les résultats MESSAGEix publiés. La différence = contribution propre de la thèse. \\

\hline

H4 &
Différenciation régionale de l'intensité des perturbations &
\textbf{12 régions IMACLIM.} Multiplier chaque \emph{demand modifier} par le ratio $r_n / r_0$ où $r_n$ est le niveau de transformation digitale de la région $n$ et $r_0$ celui de la région d'origine de l'évidence empirique. Fan \emph{et al.} fournissent les projections $r_n$ sous SSP1-5 à l'horizon 2050. &
Fan \emph{et al.} 2025 (projections EGDI, 62~pays, 12~régions, SSP1-5). &
Émissions régionales, investissement régional, emploi régional. Isoler la contribution de la fracture numérique vs.\ la stagnation économique générale. \\

\hline

Axe transversal &
Direction du changement technique~: part du travail dans la valeur ajoutée des secteurs exposés à l'IA &
\textbf{Module travail~:} paramètre de partage de la VA ($1-\pi_j$) dans le secteur services/numérique. Variante \og automation-dominante\fg{}~: $1-\pi_j$ en baisse tendancielle. Variante \og pro-worker\fg{}~: $1-\pi_j$ stable ou en hausse. Appliqué en overlay sur chaque scénario énergie. &
Séries longues de labor share sectorielle (Karabarbounis \& Neiman 2014, Autor et al.\ 2020). Taxonomie d'\textcite{acemoglu_building_2026} pour la calibration des variantes. &
Composition de la demande finale par catégorie de ménage, trajectoire d'émissions. Isoler la contribution du canal travail vs.\ le canal énergie-physique. \\

\hline
\end{tabular}
\end{table}


\subsubsection*{Questions ouvertes pour la spécification technique}

Les points suivants doivent être clarifiés avec l'équipe IMACLIM avant de fixer le protocole~:

\begin{enumerate}
\item \textbf{Créer un nouveau secteur ou modifier l'existant~?} Ajouter un secteur <<~ICT/numérique~>> à la matrice entrées-sorties (13\textsuperscript{e}~secteur) donnerait la résolution la plus fine mais exige de restructurer le tableau comptable hybride et de trouver des données physiques (TWh) et monétaires (\$) séparées pour ce secteur dans chaque région. L'alternative --- modifier les coefficients techniques du secteur <<~services et industrie légère~>> et y ajouter un sous-module dynamique --- est plus légère mais perd la distinction entre services numériques et non numériques.

\item \textbf{Comment représenter la demande des centres de données~?} Option~A~: composante additionnelle exogène de demande d'électricité dans le secteur tertiaire (analogue à ce que fait Wilson dans MESSAGEix). Option~B~: sous-module dynamique qui endogénéise la croissance de la demande en fonction du prix de l'électricité, de la productivité du calcul et du taux de marge du secteur --- plus ambitieux mais cohérent avec la logique IMACLIM.

\item \textbf{Le taux de marge sectoriel.} IMACLIM permet de différencier les taux de marge ($\pi_j$) par secteur. Les entreprises technologiques opèrent avec des marges nettes de 25 à 40\,\%, très supérieures à la moyenne industrielle. Fixer un $\pi_{services}$ différencié pour la composante numérique est techniquement simple et permet de tester si la concentration du secteur modifie les trajectoires d'investissement agrégé.

\item \textbf{La courbe de charge.} IMACLIM représente la courbe de charge électrique pour le dispatch des technologies de production. Les centres de données ont un profil de charge plat (facteur de charge $>0{,}9$), très différent du résidentiel ou de l'industrie. Ce profil modifie le merit order et peut accélérer ou freiner la pénétration des renouvelables variables selon qu'il est situé en base ou en semi-base. Cette interaction est-elle captée avec la résolution temporelle actuelle d'IMACLIM~?

\item \textbf{Le couplage IMACLIM--PyPSA.} Si la résolution temporelle d'IMACLIM ne suffit pas à capter l'interaction entre profil de charge numérique et intégration des renouvelables, un couplage soft avec un modèle de système électrique (PyPSA) pourrait compléter l'analyse. Faisabilité dans le cadre d'un doctorat~?

\item \textbf{La part du travail comme paramètre de scénario.} Le module travail d'IMACLIM-R World permet-il de faire varier le partage de la valeur ajoutée ($1-\pi_j$) comme paramètre de scénario, indépendamment de la TFP sectorielle et de l'élasticité de la wage curve~? Si oui, la direction du changement technique numérique au sens d'\textcite{acemoglu_building_2026} --- automation substitutive vs.\ création de tâches nouvelles --- peut être testée comme axe transversal des scénarios. Si non, l'axe reste un narratif qui encadre l'interprétation des résultats sans être directement paramétrisé.
\end{enumerate}


\section{Le problème de la modélisation énergie-économie}\label{sec:imaclim-probleme}

La modélisation des interactions entre systèmes énergétiques et trajectoires macroéconomiques bute sur une difficulté constitutive. Les modèles d'équilibre général calculable représentent l'économie comme un système cohérent de flux monétaires entre agents, marchés et institutions~: ils garantissent le bouclage comptable --- chaque dépense est le revenu de quelqu'un --- et permettent de calculer les effets de politiques fiscales ou commerciales sur la structure des prix, la répartition des revenus et la croissance. Mais dans un CGE standard, l'énergie est un bien parmi d'autres, substituable par des fonctions de production à élasticité constante. Une tonne de CO\textsubscript{2} n'y existe pas~; un kilowattheure non plus. La contrainte physique --- l'épuisement d'une ressource, le rendement thermodynamique d'une centrale, la courbe de charge d'un réseau électrique --- est absente du cadre, ou réduite à un paramètre d'élasticité qui condense en un scalaire ce que l'ingénieur décrit en dizaines de variables.

À l'inverse, les modèles bottom-up du système énergétique --- MARKAL, TIMES, POLES --- décrivent avec précision les technologies disponibles, leurs coûts et leurs contraintes d'intégration. Ils projettent un mix électrique compatible avec un objectif d'émissions et calculent le coût technique de l'abattement. Leur moteur macroéconomique, toutefois, est rudimentaire quand il existe~: la demande d'énergie dérive d'un PIB exogène, les effets de retour de l'investissement énergétique sur la croissance sont ignorés, et le prix n'est qu'un signal de coût partiel, déconnecté de l'équilibre des marchés des biens, du travail et du capital.

\textcite{hourcade_hybrid_2006} ont nommé cette impasse la \og fable de l'éléphant et du lapin\fg{}~: le CGE est un éléphant qui voit toute l'économie mais pas l'énergie~; le modèle bottom-up est un lapin qui voit l'énergie mais pas l'économie. Quand on met les deux dans le même ragoût, on obtient un ragoût à l'éléphant --- le cadre macroéconomique impose sa logique et les détails techniques deviennent des paramètres d'élasticité ajustables. L'alternative consiste à ne pas les mélanger dans une bouillie fonctionnelle unique mais à les faire dialoguer dans un cadre qui respecte les deux comptabilités.

Cette difficulté prend un relief particulier dès qu'on cherche à évaluer non pas le coût technique d'un objectif climatique --- exercice d'ingénierie --- mais ses conséquences macroéconomiques sur l'emploi, la distribution des revenus et le rythme de croissance. Un modèle purement technique ne peut pas voir qu'une taxe carbone, en renchérissant les coûts de production, réduit le pouvoir d'achat des ménages, contracte la demande agrégée et peut accroître le chômage --- sauf si le recyclage des recettes fiscales compense l'effet. Un modèle purement monétaire ne peut pas voir qu'une substitution charbon-renouvelable modifie les besoins en capital, le profil temporel de la courbe de charge et la géographie des investissements en réseau. Le couplage des deux dimensions n'est pas un raffinement méthodologique~; c'est une condition nécessaire pour que l'analyse ait un sens économique.

C'est ce couplage que la plateforme IMACLIM a été construite pour réaliser.

Le choix d'un cadre de modélisation pour la question posée dans cette thèse doit cependant être justifié par des critères plus précis que l'adéquation technique. \textcite{wilson_evaluating_2021}, dans une synthèse des méthodes d'évaluation des modèles intégrés d'évaluation, proposent quatre critères interdépendants qui guident ce choix. Le modèle doit être \emph{approprié} à la question scientifique posée~: sa résolution sectorielle, temporelle et géographique doit correspondre aux mécanismes qu'on cherche à tester \parencite{jakeman_ten_2006}. Ses résultats doivent être \emph{interprétables} à la lumière de sa structure et de ses hypothèses~: la devise applicable aux modèles intégrés est de communiquer des \og~insights, not numbers~\fg{} \parencite{peace_insights_2008}. Le modèle doit être \emph{crédible} auprès des communautés d'utilisateurs, ce qui suppose une confrontation aux données historiques, aux faits stylisés et aux résultats d'autres modèles. Enfin, l'analyse doit être \emph{pertinente} pour informer la compréhension scientifique et les décisions de politique climatique \parencite{cash_knowledge_2003}. Wilson et al.\ insistent sur un point décisif~: l'évaluation ne rend pas les modèles plus précis ni plus fiables dans la prédiction de l'avenir~; elle améliore leur utilité comme outils scientifiques pour comprendre les trajectoires, les options de décision et leurs conséquences. Un modèle est, selon la formule de \textcite{barlas_philosophical_1990}, non pas vrai ou faux mais situé sur un continuum d'utilité.

Ces quatre critères structurent le choix d'IMACLIM et le design expérimental du présent chapitre. L'\emph{appropriation} est assurée par la double comptabilité physique-monétaire (qui permet de suivre les térawattheures du capital computationnel en même temps que les euros investis) et par la résolution sectorielle (qui permettra d'isoler le secteur numérique). L'\emph{interprétabilité} est assurée par le balayage paramétrique systématique des hypothèses H1--H4 et par la confrontation aux résultats publiés dans MESSAGEix et TIAM. La \emph{crédibilité} est construite par la confrontation aux faits stylisés établis dans les chapitres~1 et~2, qui fournissent des contraintes que les trajectoires simulées doivent respecter. La \emph{pertinence} est garantie par le fait que la question posée~--- sous quelles conditions le couplage ICT-énergie intensifie ou atténue la consommation~--- est directement liée aux décisions de planification du système électrique que les opérateurs et les régulateurs doivent prendre dans la décennie qui vient.


\section{Genèse et architecture d'IMACLIM}\label{sec:imaclim-architecture}

IMACLIM --- pour \emph{IMpact Assessment of CLIMate policies} --- est une plateforme de modélisation développée au Centre International de Recherche sur l'Environnement et le Développement (CIRED) depuis le début des années~2000, sous l'impulsion de Jean-Charles Hourcade. Le projet naît d'un constat méthodologique~: les modèles CGE utilisés pour l'évaluation des politiques climatiques dans les années~1990, calibrés sur les matrices de comptabilité sociale de type GTAP, représentent l'énergie comme un facteur de production parmi d'autres, substituable via des fonctions CES. Cette représentation convient tant qu'on reste au voisinage de l'équilibre de référence. Elle devient inadéquate dès que l'on simule des écarts de grande amplitude --- une division par deux des émissions mondiales d'ici 2050, par exemple --- qui supposent des transformations structurelles profondes du système technique~: remplacement du parc de centrales, modification du parc automobile, rénovation du bâti, déploiement d'infrastructures nouvelles. Aucune fonction CES ne peut capturer de manière crédible de tels changements sur un siècle \parencite{hourcade_modelling_1993}.

La réponse du CIRED a consisté à construire un cadre qui sépare explicitement le court terme du long terme. Dans l'équilibre statique --- c'est-à-dire à un instant donné ---, les capacités productives et les coefficients techniques sont fixes. Les producteurs ne peuvent pas substituer le capital au travail ni le gaz au charbon d'un trimestre à l'autre~; ils ne peuvent qu'ajuster leur taux d'utilisation des capacités installées. C'est une hypothèse \emph{putty-clay}, par opposition aux fonctions de production \emph{putty-putty} des CGE standard~: la technologie est flexible \emph{ex ante} (au moment de l'investissement) mais rigide \emph{ex post} (une fois le capital installé). Cette rigidité est réaliste pour les secteurs à forte intensité capitalistique --- énergie, transport, bâtiment --- où la durée de vie des équipements se mesure en décennies \parencite{johansen_substitution_1959}.

Quand le taux d'utilisation des capacités approche la saturation, les coûts de production augmentent~: heures supplémentaires, travail de nuit, maintenance renforcée. Ce mécanisme de rendements décroissants statiques, capturé par une fonction~$\Omega$ croissante du ratio production/capacité, est la courbe d'offre inverse de chaque secteur~: elle montre comment le producteur ajuste son niveau de production en fonction des prix relatifs \parencite{corrado_capacity_1997}. Les prix de marché sont égaux au coût moyen de production augmenté d'un taux de marge sectoriel.

Entre deux équilibres statiques, des modules dynamiques sectoriels calculent l'évolution des paramètres techniques~: nouvelles capacités de production issues des investissements de la période précédente, efficacité énergétique des nouveaux équipements, mix électrique résultant des choix technologiques, motorisation du parc automobile. Ces modules sont des \og formes réduites\fg{} de modèles bottom-up --- ils incorporent l'expertise sectorielle (courbes d'apprentissage, potentiels techniques, contraintes physiques) sans la réduire à une élasticité \parencite{sassi_imaclim-r_2010}.

IMACLIM existe en deux versions principales. IMACLIM-S (statique comparatif), développé principalement par Frédéric Ghersi, compare un équilibre de référence projeté à un horizon donné avec un équilibre contrefactuel obtenu sous une politique climatique~; l'écart entre les deux fournit une mesure des effets macroéconomiques de la politique, conditionnellement aux hypothèses techniques issues d'un modèle bottom-up \parencite{ghersi_hybrid_2015}. Des versions nationales d'IMACLIM-S ont été développées pour la France \parencite{combet_fiscalite_2013, letreut_methodological_2017}, l'Arabie Saoudite \parencite{soummane_imaclim-sau_2020}, l'Inde \parencite{gupta_imaclim-ind_2018}, le Brésil \parencite{lefevre_hybridization_2016} et l'Afrique du Sud \parencite{schers_green_2015}, via une plateforme générique IMACLIM-Country\footnote{La plateforme IMACLIM-Country, principalement développée par Gaëlle Le~Treut avec les contributions d'Emmanuel Combet, Ruben Bibas et Julien Lefèvre, est disponible sous licence CeCILL \parencite{letreut_imaclim-country_2019}. Le jeu de données hybrides pour la France, année de base 2010, est disponible sur Mendeley Data \parencite{letreut_hybrid_2018}.}. IMACLIM-R (récursif dynamique), dont le développement a impliqué Olivier Sassi, Renaud Crassous, Henri Waisman, Céline Guivarch et Ruben Bibas, enchaîne les équilibres statiques d'année en année, avec des modules dynamiques entre chaque pas de temps \parencite{bibas_imaclim-r_2016}. Les anticipations des agents sont myopiques~: ils décident sur la base des signaux de prix courants et passés, sans connaissance du futur. L'investissement n'est pas le résultat d'une optimisation intertemporelle~; il est alloué entre secteurs et régions selon la rentabilité anticipée, telle que perçue par des agents qui peuvent se tromper.

Le trait distinctif d'IMACLIM, commun aux deux versions, est la double comptabilité physique et monétaire. Chaque flux économique a un miroir physique. Les secteurs énergétiques --- pétrole, gaz, charbon, raffinage, électricité --- sont comptabilisés simultanément en unités physiques (ktep, TWh) et en unités monétaires. Le lien entre les deux est assuré par un vecteur de prix spécifiques par agent --- les \og marges spécifiques\fg{} --- qui reflète le fait qu'un même bien énergétique n'est pas vendu au même prix à un industriel, à un ménage et à un exportateur. Cette hybridation, qui repose techniquement sur la réconciliation des tableaux entrée-sortie de la comptabilité nationale avec les bilans énergétiques de l'AIE, a des conséquences analytiques profondes~: elle permet de suivre les intensités énergétiques sectorielles en unités physiques, de calculer les émissions de CO\textsubscript{2} directement à partir des consommations physiques de combustibles fossiles, et de voir comment une politique de prix modifie simultanément les flux monétaires et les flux physiques --- sans réduire l'un à l'autre \parencite{combet_construction_2014}.


\section{Le second-best~: ce qu'IMACLIM fait que les autres ne font pas}\label{sec:imaclim-secondbest}

La plupart des modèles intégrés utilisés dans les exercices d'évaluation du GIEC --- REMIND, WITCH, MESSAGE, GCAM --- reposent sur une hypothèse de premier-best~: l'économie est à l'optimum, les marchés sont parfaits, les facteurs de production sont pleinement utilisés, et les agents anticipent correctement le futur. Dans ce cadre, le coût d'une politique climatique se mesure comme l'écart entre deux optima --- celui avec contrainte climatique et celui sans ---, et cet écart est nécessairement un coût. Le double dividende (une politique climatique qui améliore le bien-être net, co-bénéfices inclus) est impossible par construction~: on ne peut pas améliorer un optimum en ajoutant une contrainte.

IMACLIM part d'une prémisse différente. L'économie de référence n'est pas à l'optimum. Trois types d'imperfections sont représentés explicitement.

Le marché du travail ne s'ajuste pas par les prix. IMACLIM utilise une \emph{wage curve} --- une relation empirique négative entre le salaire réel et le taux de chômage --- à la place de la courbe d'offre de travail flexible des CGE standard \parencite{blanchflower_introduction_1995}. Le chômage est endogène~: si une taxe carbone augmente les coûts de production, réduit la demande et contracte l'emploi, le taux de chômage monte~; il ne disparaît pas par une baisse instantanée des salaires. \textcite{guivarch_costs_2011} ont montré que l'élasticité de la wage curve est un paramètre déterminant des coûts de stabilisation~: avec des marchés du travail très flexibles (élasticité de $-7$), les résultats d'IMACLIM convergent vers la fourchette GIEC, soit moins de 2\,\% de perte de PIB en 2030 pour un objectif de 550\,ppm éq-CO\textsubscript{2}~; avec des valeurs empiriquement estimées (élasticité de $-0{,}1$), les coûts augmentent d'un facteur cinq à dix. Ce résultat traduit le fait que la transition énergétique est aussi un problème de marché du travail, et que des politiques d'accompagnement --- recyclage des recettes carbone en subventions à l'emploi --- sont nécessaires pour contenir les coûts.

Les anticipations sont imparfaites. Dans IMACLIM-R, les agents ne connaissent pas les prix futurs du pétrole, de l'électricité ni du carbone. Ils extrapolent les tendances passées. Leurs décisions d'investissement peuvent se révéler inadaptées \emph{ex post}, ce qui crée des capacités sous-utilisées ou des pénuries. Cette myopie est particulièrement conséquente dans les secteurs à forte inertie capitalistique~: un investissement dans une centrale à charbon en 2025, construit sur l'anticipation d'un prix du carbone nul, devient un actif échoué si un prix du carbone est introduit en 2030. \textcite{waisman_imaclim-r_2012} ont montré que cette interaction entre inertie technique et anticipations imparfaites est le principal déterminant des coûts de transition à court terme --- des coûts que les modèles à anticipations parfaites, par construction, ne voient pas.

Le capital installé est sectoriel et non mobile. Une usine sidérurgique ne se reconvertit pas en ferme solaire. Dans l'équilibre statique d'IMACLIM, les capacités productives de chaque secteur sont fixes~; seul le taux d'utilisation varie. Quand les prix relatifs changent sous l'effet d'une taxe carbone, les secteurs dont les capacités deviennent non compétitives ne libèrent pas instantanément leur capital pour d'autres usages. Le capital reste immobilisé, les travailleurs ne se réallouent pas sans friction, et l'économie traverse une phase de sous-utilisation des facteurs de production. Ce mécanisme est ce qui produit les \og coûts de transition\fg{} observés dans les simulations IMACLIM~: des pertes de PIB concentrées sur les deux premières décennies de la politique climatique, suivies d'un rattrapage partiel à mesure que le stock de capital se renouvelle \parencite{crassous_endogenous_2006}.

La conjonction de ces trois imperfections fait d'IMACLIM un modèle de second-best. Les conséquences sont symétriques~: la politique climatique peut aggraver les distorsions existantes (si elle augmente le chômage) ou les corriger (si elle réduit la dépendance au pétrole et les rentes fossiles). Le double dividende n'est pas exclu. \textcite{rozenberg_climate_2010} ont chiffré le co-bénéfice des politiques climatiques comme couverture contre l'incertitude pétrolière à 11\,500 milliards de dollars en valeur actualisée nette. \textcite{combet_carbon_2010} ont montré que l'impact distributif d'une taxe carbone dépend du schéma de recyclage des recettes fiscales au moins autant que du niveau de la taxe~: trois schémas de recyclage testés sur la France~2004 produisent des effets opposés en termes d'emploi et d'équité, pour une même trajectoire d'émissions.


\section{IMACLIM dans le paysage des modèles intégrés}\label{sec:imaclim-paysage}

Il serait malhonnête de présenter IMACLIM sans situer ce qu'il ne fait pas. Chaque IAM incarne un compromis entre la richesse de la représentation technique, la cohérence macroéconomique et la tractabilité computationnelle. IMACLIM privilégie les deux premiers au prix de concessions sur le troisième.

REMIND \parencite{bauer_remind-r_2012} et WITCH \parencite{bosetti_witch_2006} sont des modèles d'optimisation intertemporelle à anticipations parfaites. Ils calculent la trajectoire qui maximise le bien-être actualisé sous contrainte climatique. Leur force est la cohérence dynamique~: chaque décision d'investissement est prise en connaissance du futur, ce qui élimine les actifs échoués par construction et produit des profils de coûts lisses, croissant régulièrement dans le temps. Leur faiblesse est l'hypothèse même~: les agents réels n'optimisent pas sur un siècle, et les coûts de transition --- les pertes concentrées sur les premières décennies --- sont absents des résultats. La comparaison RECIPE \parencite{luderer_regional_2012} a montré que REMIND et WITCH produisent des prix du carbone modérés et des coûts régionaux peu dispersés, tandis qu'IMACLIM génère des prix du carbone plus élevés à court terme, avec des coûts beaucoup plus sensibles au schéma de répartition des permis. L'effet multiplicateur macroéconomique des transferts carbone atteint un coefficient de~2,8 pour la Chine dans IMACLIM --- un effet absent des modèles d'optimisation.

GCAM (Pacific Northwest National Laboratory) et MESSAGE (IIASA) sont des modèles à structure modulaire avec une résolution sectorielle et géographique plus fine qu'IMACLIM --- GCAM distingue 32~régions et représente l'usage des sols de manière endogène. Leur représentation macroéconomique est toutefois limitée~: la demande d'énergie dérive de projections de PIB et de population exogènes, et les boucles de rétroaction entre prix de l'énergie, emploi et croissance sont soit absentes, soit simplifiées. IMAGE (PBL) se distingue par son traitement détaillé de la biosphère et des cycles biogéochimiques~; ce n'est pas la question que nous posons.

Le positionnement d'IMACLIM dans ce paysage peut se résumer en trois propositions. Premièrement, IMACLIM est le seul IAM largement diffusé qui représente simultanément le sous-emploi endogène, les anticipations myopiques et la double comptabilité physique/monétaire. Deuxièmement, il est le seul à produire des profils temporels de coûts non monotones --- avec des coûts de transition élevés à court terme et un rattrapage potentiel à long terme ---, ce qui en fait un outil adapté à l'analyse des trajectoires plutôt que des états finaux. Troisièmement, il n'optimise pas, ne produit pas de \og trajectoires optimales\fg{}, ne représente pas l'usage des sols de manière endogène, et sa résolution régionale (12~régions dans IMACLIM-R World) est moins fine que celle de GCAM ou de MESSAGE.


Une évaluation systématique récente de dix modèles représentatifs --- REMIND, MESSAGE-MACRO, GEM-E3, GCAM, IMACLIM-R, DICE, GEMMES, E3ME, MEDEAS et EUROGREEN --- confirme ce positionnement \parencite{arias_modeling_2026}\footnote{Julien Lefèvre (CIRED, équipe IMACLIM) est co-auteur.}. Les auteurs identifient cinq catégories de shortcomings~: paradigme d'équilibre~(3.1), représentation de la finance et du crédit~(3.2), énergie et efficacité énergétique~(3.3), lien énergie-économie via fonctions de production agrégées~(3.4), et matériaux~(3.5). IMACLIM-R échappe aux shortcomings liés aux fonctions de production agrégées~--- Cambridge Capital Controversy~(3.4.1), incertitude des paramètres CES~(3.4.2), sensibilité insuffisante du PIB à l'énergie~(3.4.3)~--- précisément parce que le PIB y est calculé comme somme de la production sectorielle via les tables I-O, et non via une APF de type CES ou Cobb-Douglas. Le mark-up pricing, le chômage involontaire et la myopie des agents évitent les critiques relatives à la concurrence parfaite~(3.1.3) et aux anticipations parfaites~(3.1.1). Les shortcomings qui s'appliquent à IMACLIM-R concernent la productivité du travail et l'intensité énergétique exogènes~(3.4.4), l'absence de crédit endogène~(3.2.1) et de comptabilité stock-flow consistent~(3.2.2), la captation partielle des effets rebond~(3.3.1), et l'absence d'EROI dynamique~(3.3.2). Le papier conclut que \og using input-output tables with multiple sectors and myopic decision-making with recursive dynamics helps overcome some shortcomings associated with equilibrium-based modeling\fg{}, tout en notant qu'aucun des dix modèles évalués ne performe suffisamment bien sur l'ensemble des critères.

\section{Pourquoi IMACLIM pour la question du numérique}\label{sec:imaclim-numerique}

Le choix d'IMACLIM comme cadre d'accueil pour l'endogénéisation du couplage numérique-énergie repose sur trois arguments structurels.

Le premier est architectural. Dans IMACLIM, les coefficients techniques --- les quantités physiques d'énergie, de travail et d'intrants intermédiaires nécessaires pour produire une unité de bien --- sont des paramètres de l'équilibre statique, mis à jour entre deux équilibres par des modules dynamiques sectoriels. Ajouter un module \og infrastructure numérique\fg{} qui calcule, en fonction des investissements dans les centres de données~($K_{DC}$), des contrats d'approvisionnement électrique~($P_{PPA}$) et de l'intensité énergétique du calcul~($\eta_{AI}$), les nouveaux coefficients input-output pour les secteurs des services et de l'industrie revient à suivre le même protocole que celui utilisé pour le module électricité ou le module transport. Il n'y a pas d'adaptation \emph{ad hoc} du cadre~; c'est son mode de fonctionnement normal.

Le deuxième est comptable. L'infrastructure numérique est un consommateur d'électricité dont la demande physique --- en TWh --- est aujourd'hui du même ordre de grandeur que celle de secteurs industriels entiers, et dont la trajectoire future dépend de paramètres techniques (efficacité des GPU, PUE des centres de données, taux d'utilisation des serveurs) qui ne peuvent pas être résumés par une élasticité de substitution. La double comptabilité physique/monétaire d'IMACLIM permet de suivre simultanément les térawattheures consommés par le capital computationnel et les euros investis dans ce capital --- et de voir comment l'un affecte l'autre via le vecteur de prix. Dans un CGE purement monétaire, un centre de données qui consomme 100\,MW est indistinguable d'une usine qui consomme 100\,MW~; dans IMACLIM, les coefficients techniques les distinguent.

Le troisième est macroéconomique. Le numérique modifie la productivité du travail, la structure sectorielle de l'emploi et les schémas de consommation des ménages. Ces effets passent par le marché du travail --- exactement le canal qu'IMACLIM, grâce à sa wage curve et à son chômage endogène, est conçu pour capturer. Un modèle de premier-best, dans lequel le plein emploi est garanti par hypothèse, ne peut pas voir qu'une accélération de l'automatisation par l'IA accroît la productivité agrégée tout en détruisant des emplois dans les secteurs routinisables --- ni que l'effet net sur le PIB dépend de la rigidité du marché du travail et de la politique de redistribution. IMACLIM le peut. Plus précisément, \textcite{acemoglu_building_2026} montrent que même au sein des technologies dites complémentaires, les effets sur la part du travail divergent~: une technologie \emph{labor-augmenting} la laisse inchangée, une technologie créatrice de tâches nouvelles l'augmente. Tester cette distinction exige un modèle où la distribution fonctionnelle des revenus affecte la composition de la demande, ce qui est le cas d'IMACLIM grâce à la wage curve, au chômage endogène et à la différenciation des consommations par catégorie de ménage. C'est une propriété nécessaire pour une analyse qui ne se contente pas de projeter la demande électrique des centres de données, mais qui veut comprendre comment le déploiement du capital computationnel interagit avec la trajectoire macroéconomique d'une transition bas-carbone.
